\section{AHSME 1967}
        \begin{enumerate}
        	\item (\textit{AHSME 1967}) The three-digit number $2a3$ is added to the number $326$ to give the three-digit number $5b9$.  If $5b9$ is divisible by 9, then $a+b$ equals


 $\text{(A)}\ 2\qquad\text{(B)}\ 4\qquad\text{(C)}\ 6\qquad\text{(D)}\ 8\qquad\text{(E)}\ 9$

 \textbf{Answer: }C



\textbf{Solution}

If $5b9$ is divisible by $9$, this must mean that $5 + b + 9$ is a multiple of $9$. So, 
 \[5 + b + 9 = 9, 18, 27, 36...\] 

 Because $5 + 9 = 14$ and $b$ is in between 0 and 9, 

 
 \[5 + b + 9 = 18\]
 \[b = 4\]

 The question states that
 \[2a3 + 326 = 549\]so 
 \[2a3 = 549 - 326\]
 \[2a3 = 223\]
 \[a = 2\]

 \[a + b = 6\]which is answer choice $\boxed{C}$.

 


	\item (\textit{AHSME 1967}) An equivalent of the expression


 $\left(\frac{x^2+1}{x}\right)\left(\frac{y^2+1}{y}\right)+\left(\frac{x^2-1}{y}\right)\left(\frac{y^2-1}{x}\right)$, $xy \not= 0$,


 is:


 $\text{(A)}\ 1\qquad\text{(B)}\ 2xy\qquad\text{(C)}\ 2x^2y^2+2\qquad\text{(D)}\ 2xy+\frac{2}{xy}\qquad\text{(E)}\ \frac{2x}{y}+\frac{2y}{x}$

 \textbf{Answer: }D



\textbf{Solution 1}


 \[\left(\frac{x^2+1}{x}\right)\left(\frac{y^2+1}{y}\right)+\left(\frac{x^2-1}{y}\right)\left(\frac{y^2-1}{x}\right)\]
 \[\frac{(x^2+1)(y^2+1)}{xy} + \frac{(x^2-1)(y^2-1)}{xy}\]
 \[\frac{(x^2y^2+x^2+y^2+1) + (x^2y^2-x^2-y^2+1)}{xy}\]
 \[\frac{x^2y^2+x^2+y^2+1 + x^2y^2-x^2-y^2+1}{xy}\]
 \[\frac{2x^2y^2+2}{xy}\]
 \[\frac{2x^2y^2}{xy} + \frac{2}{xy}\]
 \[\frac{2xy * xy}{xy} + \frac{2}{xy}\]
 \[2xy + \frac{2}{xy}\]The answer is $\boxed{D}$




\textbf{Solution 2 (Answer Choices)}

The denominator before simplification is $xy$. As the constant term will not be divisible by the denominator, we inspect that the fraction with only the constant term in the numerator is $\frac{(1 \cdot 1) + (-1 \cdot -1)}{xy}$. The only answer choice that matches this is $\boxed{D}$.~megaboy6679

 


 


	\item (\textit{AHSME 1967}) The side of an equilateral triangle is $s$.  A circle is inscribed in the triangle and a square is inscribed in the circle.  The area of the square is:


 $\text{(A)}\ \frac{s^2}{24}\qquad\text{(B)}\ \frac{s^2}{6}\qquad\text{(C)}\ \frac{s^2\sqrt{2}}{6}\qquad\text{(D)}\ \frac{s^2\sqrt{3}}{6}\qquad\text{(E)}\ \frac{s^2}{3}$

 \textbf{Answer: }B



\textbf{Solution}

The radius of our circle is $\frac{\sqrt{3}}{6}s$. This means the square has side length $\frac{\sqrt{6}}{6}s$ which has area of $\frac{s^2}{6}=\fbox{B}$.

 


	\item (\textit{AHSME 1967}) Given $\frac{\log{a}}{p}=\frac{\log{b}}{q}=\frac{\log{c}}{r}=\log{x}$, all logarithms to the same base and $x \not= 1$.  If $\frac{b^2}{ac}=x^y$, then $y$ is:


 $\text{(A)}\ \frac{q^2}{p+r}\qquad\text{(B)}\ \frac{p+r}{2q}\qquad\text{(C)}\ 2q-p-r\qquad\text{(D)}\ 2q-pr\qquad\text{(E)}\ q^2-pr$

 \textbf{Answer: }C



\textbf{Solution}

We are given:
 \[\frac{b^2}{ac} = x^y\]

 Taking the logarithm on both sides:
 \[\log{\left(\frac{b^2}{ac}\right)} = \log{x^y}\]

 Using the properties of logarithms:
 \[2\log{b} - \log{a} - \log{c} = y \log{x}\]

 Substituting the values given in the problem statement:
 \[2q \log{x} - p \log{x} - r \log{x} = y \log{x}\]

 Since $x \neq 1$, dividing each side by $\log{x}$ we get:
 \[y = \boxed{\textbf{(C) } 2q - p - r}\]



	\item (\textit{AHSME 1967}) A triangle is circumscribed about a circle of radius $r$ inches.  If the perimeter of the triangle is $P$ inches and the area is $K$ square inches, then $\frac{P}{K}$ is:


 $\text{(A)}\text{independent of the value of} \; r\qquad\text{(B)}\ \frac{\sqrt{2}}{r}\qquad\text{(C)}\ \frac{2}{\sqrt{r}}\qquad\text{(D)}\ \frac{2}{r}\qquad\text{(E)}\ \frac{r}{2}$

 \textbf{Answer: }D



\textbf{Solution}

The area $K$ of the triangle can be expressed in terms of its inradius $r$ and its semiperimeter $s$ as:
 \[K = r \times s = r \times \frac{P}{2}\]

 So, $\frac{P}{K} = \boxed{\textbf{(C) } \frac{2}{r}}$.

 


	\item (\textit{AHSME 1967}) If $f(x)=4^x$ then $f(x+1)-f(x)$ equals:


 $\text{(A)}\ 4\qquad\text{(B)}\ f(x)\qquad\text{(C)}\ 2f(x)\qquad\text{(D)}\ 3f(x)\qquad\text{(E)}\ 4f(x)$

 \textbf{Answer: }D



\textbf{Solution}

The desired expression is equal to $4^{x+1} - 4^{x}$Using the fact that $4^{x+1}$=$4^{x}*4$, we see that the answer is$3*4^{x}$$\fbox{D}$



	\item (\textit{AHSME 1967}) If $\frac{a}{b}<-\frac{c}{d}$ where $a$, $b$, $c$, and $d$ are real numbers and $bd \not= 0$, then:


 $\text{(A)}\ a \; \text{must be negative} \qquad \text{(B)}\ a \; \text{must be positive} \qquad$
$\text{(C)}\ a \; \text{must not be zero} \qquad \text{(D)}\ a \; \text{can be negative or zero, but not positive } \\ \text{(E)}\ a \; \text{can be positive, negative, or zero}$

 \textbf{Answer: }E



\textbf{Solution}

Notice that $b$ and $d$ are irrelevant to the problem.  Also notice that the negative sign is irrelevant since $c$ can be any real number.  So the question is equivalent to asking what values $a$ can take on given the inequality $a<c$.  The answer, is of course, that $a$ could be anything, or $\boxed{\textbf{(E) } \ a \; \text{can be positive, negative, or zero}}$



	\item (\textit{AHSME 1967}) To $m$ ounces of a $m\%$ solution of acid, $x$ ounces of water are added to yield a $(m-10)\%$ solution.  If $m>25$, then $x$ is 


 $\textbf{(A)}\ \frac{10m}{m-10} \qquad \textbf{(B)}\ \frac{5m}{m-10} \qquad \textbf{(C)}\ \frac{m}{m-10} \qquad \textbf{(D)}\ \frac{5m}{m-20} \\ \textbf{(E)}\ \text{not determined by the given information}$

 \textbf{Answer: }A



\textbf{Solution}

$\fbox{A}$



	\item (\textit{AHSME 1967}) Let $K$, in square units, be the area of a trapezoid such that the shorter base, the altitude, and the longer base, in that order, are in arithmetic progression.  Then:


 $\textbf{(A)}\ K \; \text{must be an integer} \qquad \textbf{(B)}\ K \; \text{must be a rational fraction} \\ \textbf{(C)}\ K \; \text{must be an irrational number} \qquad \textbf{(D)}\ K\;  \text{must be an integer or a rational fraction} \qquad$
$\textbf{(E)}\ \text{taken alone neither} \; \textbf{(A)} \; \text{nor} \; \textbf{(B)} \; \text{nor} \; \textbf{(C)} \; \text{nor} \; \textbf{(D)} \; \text{is true}$

 \textbf{Answer: }E



\textbf{Solution}

From the problem we can set the altitude equal to $a$, the shorter base equal to $a-d$, and the longer base equal to $a+d$. By the formula for the area of a trapezoid, we have $K=a^2$. However, since $a$ can equal any real number $(3, 2.7, \pi)$, none of the statements $A, B, C, D$ need to be true, so the answer is $\fbox{E}$.

 


	\item (\textit{AHSME 1967}) If $\frac{a}{10^x-1}+\frac{b}{10^x+2}=\frac{2 \cdot 10^x+3}{(10^x-1)(10^x+2)}$ is an identity for positive rational values of $x$, then the value of $a-b$ is:


 $\textbf{(A)}\ \frac{4}{3} \qquad \textbf{(B)}\ \frac{5}{3} \qquad \textbf{(C)}\ 2 \qquad \textbf{(D)}\ \frac{11}{4} \qquad \textbf{(E)}\ 3$

 \textbf{Answer: }A



\textbf{Solution}

Given the equation:
 \[\frac{a}{10^x-1}+\frac{b}{10^x+2}=\frac{2 \cdot 10^x+3}{(10^x-1)(10^x+2)}\]

 Let's simplify by letting $y = 10^x$. The equation becomes:
 \[\frac{a}{y-1}+\frac{b}{y+2}=\frac{2y+3}{(y-1)(y+2)}\]

 Multiplying each term by the common denominator $(y-1)(y+2)$, we obtain:
 \[a(y+2) + b(y-1) = 2y+3\]

 For the equation to be an identity, the coefficients of like terms on both sides must be equal. Equating the coefficients of $y$ and the constant terms, we get the system of equations:

 
 \begin{align*} a+b &= 2 \quad \text{(from coefficients of } y \text{)} \\ 2a - b &= 3 \quad \text{(from constant terms)} \end{align*}


 Solving this system, we find:
 \[a = \frac{5}{3} \quad \text{and} \quad b = \frac{1}{3}\]

 Thus, the difference is:
 \[a-b = \boxed{\textbf{(A) } \frac{4}{3}}\]



	\item (\textit{AHSME 1967}) If the perimeter of rectangle $ABCD$ is $20$ inches, the least value of diagonal $\overline{AC}$, in inches, is:


 $\textbf{(A)}\ 0\qquad \textbf{(B)}\ \sqrt{50}\qquad \textbf{(C)}\ 10\qquad \textbf{(D)}\ \sqrt{200}\qquad \textbf{(E)}\ \text{none of these}$

 \textbf{Answer: }B



\textbf{Solution}

For rectangle $ABCD$ with perimeter 20, the diagonal $AC$ is given by:
 \[AC = \sqrt{l^2 + w^2}\]To minimize $AC$, $l$ and $w$ should be equal (i.e., the rectangle is a square). Thus, $l = w = 5$.So, the minimum $AC$ is:
 \[AC = \sqrt{5^2 + 5^2} = \boxed{\textbf{(B) } \sqrt{50}}\]~ proloto

 


	\item (\textit{AHSME 1967}) If the (convex) area bounded by the x-axis and the lines $y=mx+4$, $x=1$, and $x=4$ is $7$, then $m$ equals:


 $\textbf{(A)}\ -\frac{1}{2}\qquad \textbf{(B)}\ -\frac{2}{3}\qquad \textbf{(C)}\ -\frac{3}{2} \qquad \textbf{(D)}\ -2 \qquad \textbf{(E)}\ \text{none of these}$

 \textbf{Answer: }B


	\item (\textit{AHSME 1967}) A triangle $ABC$ is to be constructed given a side $a$ (opposite angle $A$). angle $B$, and $h_c$, the altitude from $C$.  If $N$ is the number of noncongruent solutions, then $N$
$\textbf{(A)}\ \text{is} \; 1\qquad \textbf{(B)}\ \text{is} \; 2\qquad \textbf{(C)}\ \text{must be zero}\qquad \textbf{(D)}\ \text{must be infinite}\qquad \textbf{(E)}\ \text{must be zero or infinite}$

 \textbf{Answer: }E


	\item (\textit{AHSME 1967}) Let $f(t)=\frac{t}{1-t}$, $t \not= 1$.  If $y=f(x)$, then $x$ can be expressed as 


 $\textbf{(A)}\ f\left(\frac{1}{y}\right)\qquad \textbf{(B)}\ -f(y)\qquad \textbf{(C)}\ -f(-y)\qquad \textbf{(D)}\ f(-y)\qquad \textbf{(E)}\ f(y)$

 \textbf{Answer: }C



\textbf{Solution}

Since we know that $y=f(x)$, we can solve for $y$ in terms of $x$. This gives us

 $y=\frac{x}{1-x}$
$\Rightarrow y(1-x)=x$
$\Rightarrow y-yx=x$
$\Rightarrow y=yx+x$
$\Rightarrow y=x(y+1)$
$\Rightarrow x=\frac{y}{y+1}$

 Therefore, we want to find the function with $y$ that outputs $\frac{y}{y+1}$ Listing out the possible outputs from each of the given functions we get$f\left(\frac{1}{y}\right)=\frac{1}{y-1}$
$f(y)=\frac{1}{1-y}$
$f(-y)=\frac{-y}{y+1}$
$-f(y)=\frac{y}{y-1}$
$-f(-y)=\frac{y}{y+1}$

 Since $-f(-y)=\frac{y}{y+1}=x$ the answer must be $\boxed{C}$.

 


	\item (\textit{AHSME 1967}) The difference in the areas of two similar triangles is $18$ square feet, and the ratio of the larger area to the smaller is the square of an integer.  The area of the smaller triangle, in square feet, is an integer, and one of its sides is $3$ feet.  The corresponding side of the larger triangle, in feet, is:


 $\textbf{(A)}\ 12\quad \textbf{(B)}\ 9\qquad \textbf{(C)}\ 6\sqrt{2}\qquad \textbf{(D)}\ 6\qquad \textbf{(E)}\ 3\sqrt{2}$

 \textbf{Answer: }D



\textbf{Solution}

If the areas of the two triangles are $A_1$ and $A_2$ with $A_1 > A_2$, we are given that $A_1 - A_2 = 18$ and $\frac{A_1}{A_2} = k^2$

 Plugging in the second equation into the first leads to $(k^2 - 1)A_2 = 18$.  

 If $A_2$ is an integer, it can only be a factor of $18$ - namely $1, 2, 3, 6, 9, 18$.  

 Since $A_1 = A_2 + 18$, this would lead to $A_1 = 19, 20, 21, 24, 27, 36$.  

 This would lead to $\frac{A_1}{A_2} = 19, 10, 7, 4, 3, 2$, and the only square is $4$.  

 Thus, the ratio of the areas is $4$, and hence the ratio of the sides is $\sqrt{4} = 2$.  The corresponding side has a length of $3 \cdot 2 = 6$, which is option $\fbox{D}$.

 


	\item (\textit{AHSME 1967}) Let the product $(12)(15)(16)$, each factor written in base $b$, equal $3146$ in base $b$.  Let $s=12+15+16$, each term expressed in base $b$.  Then $s$, in base $b$, is 


 $\textbf{(A)}\ 43\qquad \textbf{(B)}\ 44\qquad \textbf{(C)}\ 45\qquad \textbf{(D)}\ 46\qquad \textbf{(E)}\ 47$

 \textbf{Answer: }B



\textbf{Solution}

Converting everything into base $10$, we have $(b + 2)(b+5)(b+6) = 3b^3 + b^2 + 4b + 6$.  Looking ahead, the constant term of the polynomial will be $2*5*6 -6 = 54$. By the Rational Root Theorem, the only possible integer roots are $1, 2, 3, 6, 9, 18, 27, 54$.  Bases $1, 2, 3, 6$ do not have a $6$ as a digit.  Testing $b=9$ gives a solution that works.

 Therefore, we are working in base $9$.  Adding the units place in base $9$, $2 + 5 + 6 = 14_9$, so we carry the $1$ to get a total of $44_9$, which is option $\fbox{B}$.

 


	\item (\textit{AHSME 1967}) If $r_1$ and $r_2$ are the distinct real roots of $x^2+px+8=0$, then it must follow that:


 $\textbf{(A)}\ |r_1+r_2|>4\sqrt{2}\qquad \textbf{(B)}\ |r_1|>3 \; \text{or} \; |r_2| >3 \\ \textbf{(C)}\ |r_1|>2 \; \text{and} \; |r_2|>2\qquad \textbf{(D)}\ r_1<0 \; \text{and} \; r_2<0\qquad \textbf{(E)}\ |r_1+r_2|<4\sqrt{2}$

 \textbf{Answer: }A



\textbf{Solution}

We are given that the roots are real, so the discriminant is positive, which means $p^2 - 4(8)(1) > 0$.  This leads to $|p| > 4\sqrt{2}$.  By Vieta, the sum of the roots is $-p$, so we have $|-(r_1 + r_2)| \ge 4\sqrt{2}$, or $|r_1 + r_2| > 4\sqrt{2}$, which is option $\fbox{A}$.

 


	\item (\textit{AHSME 1967}) If $x^2-5x+6<0$ and $P=x^2+5x+6$ then


 $\textbf{(A)}\ P \; \text{can take any real value} \qquad \textbf{(B)}\ 20<P<30\\ \textbf{(C)}\ 0<P<20 \qquad \textbf{(D)}\ P<0 \qquad \textbf{(E)}\ P>30$

 \textbf{Answer: }B



\textbf{Solution}

We are given that $x^2 - 5x + 6 < 0$, which, when factored, gives $(x - 2)(x-3) < 0$.  This has a solution of $2<x<3$, because the original quadratic is $\cup$-shaped, and thus dips below the x-axis between the roots.

 Since $x^2 + 5x + 6$ has a vertex minimum at $x = -\frac{5}{2}$, so it is increasing on the interval $[2, 3]$.  Thus, evaluating $P$ at $x=2$ and $x=3$ will give our bounds, and doing so gives $20 < P < 30$, or $\fbox{B}$.

 


	\item (\textit{AHSME 1967}) The area of a rectangle remains unchanged when it is made $2 \frac{1}{2}$ inches longer and $\frac{2}{3}$ inch narrower, or when it is made $2 \frac{1}{2}$ inches shorter and $\frac{4}{3}$ inch wider.  Its area, in square inches, is:


 $\textbf{(A)}\ 30\qquad \textbf{(B)}\ \frac{80}{3}\qquad \textbf{(C)}\ 24\qquad \textbf{(D)}\ \frac{45}{2}\qquad \textbf{(E)}\ 20$

 \textbf{Answer: }E



\textbf{Solution}

We are given $xy = (x+\frac{5}{2})(y-\frac{2}{3}) = (x - \frac{5}{2})(y + \frac{4}{3})$

 FOILing each side gives:

 $xy = xy - \frac{2}{3}x + \frac{5}{2}y - \frac{5}{3} = xy + \frac{4}{3}x - \frac{5}{2}y - \frac{10}{3}$

 Taking the last two parts and moving everything to the left gives:

 $-2x + 5y + \frac{5}{3} = 0$

 Taking the first two parts and multiplying by $3$ gives $-2x + \frac{15}{2}y - 5 = 0$

 Solving both equations for $-2x$ and setting them equal to each other gives $5y + \frac{5}{3} = \frac{15}{2}y - 5$, which leads to $y = \frac{8}{3}$

 Plugging that in to $-2x + 5y + \frac{5}{3} = 0$ gives $x = \frac{15}{2}$.

 The area of the rectangle is $xy = 20$, or $\fbox{E}$.

 


	\item (\textit{AHSME 1967}) A circle is inscribed in a square of side $m$, then a square is inscribed in that  circle, then a circle is inscribed in the latter square, and so on.  If $S_n$ is the sum of the areas of the first $n$ circles so inscribed, then, as $n$ grows beyond all bounds, $S_n$ approaches:


 $\textbf{(A)}\ \frac{\pi m^2}{2}\qquad \textbf{(B)}\ \frac{3\pi m^2}{8}\qquad \textbf{(C)}\ \frac{\pi m^2}{3}\qquad \textbf{(D)}\ \frac{\pi m^2}{4}\qquad \textbf{(E)}\ \frac{\pi m^2}{8}$

 \textbf{Answer: }A



\textbf{Solution}

All answers correctly grow as $m^2$, so we let $m=1$.

 The radius of the first circle is $\frac{1}{2}$, so its area is $\frac{\pi}{4}$.

 The diagonal of the second square is the diameter of the first circle, which is $1$.  Therefore, the side length of the square is $\frac{\sqrt{2}}{2}$.

 Now we note that the picture is self-similar; if we erase the outer square, erase the outer circle, rotate the picture, and dilate the picture from the center in both the x- and y-directions by an equal scaling factor, we will get the original picture.  Therefore, the side lengths of successive squares form a geometric sequence with common ratio $R$, as do the radii of the circles.  On the other hand, the areas of the squares (and areas of the circles) form a geometric sequence with common ratio $R^2$.

 Since the first square has side $1$ and the second square has side $\frac{\sqrt{2}}{2}$, we know $R = \frac{\sqrt{2}}{2}$, and $R^2 = \frac{1}{2}$.  

 Since the area of the first circle is $\frac{\pi}{4}$, and the common ratio of areas is $\frac{1}{2}$, the sum of all the areas of the circles is the sum of the infinite geometric sequence, which is $\frac{A_1}{1 - r} = \frac{\frac{\pi}{4}}{1 - \frac{1}{2}} = \frac{\pi}{2}$.  This is for side length $m=1$, and as noted before, there should be an $m^2$ factor in addition to this number to generalize it from the unit square.  This gives answer $\fbox{A}$.

 


	\item (\textit{AHSME 1967}) In right triangle $ABC$ the hypotenuse $\overline{AB}=5$ and leg $\overline{AC}=3$.  The bisector of angle $A$ meets the opposite side in $A_1$.  A second right triangle $PQR$ is then constructed with hypotenuse $\overline{PQ}=A_1B$ and leg $\overline{PR}=A_1C$.  If the bisector of angle $P$ meets the opposite side in $P_1$, the length of $PP_1$ is:


 $\textbf{(A)}\ \frac{3\sqrt{6}}{4}\qquad \textbf{(B)}\ \frac{3\sqrt{5}}{4}\qquad \textbf{(C)}\ \frac{3\sqrt{3}}{4}\qquad \textbf{(D)}\ \frac{3\sqrt{2}}{2}\qquad \textbf{(E)}\ \frac{15\sqrt{2}}{16}$

 \textbf{Answer: }B



\textbf{Solution}

$\fbox{B}$



	\item (\textit{AHSME 1967}) For natural numbers, when $P$ is divided by $D$, the quotient is $Q$ and the remainder is $R$.  When $Q$ is divided by $D'$, the quotient is $Q'$ and the remainder is $R'$.  Then, when $P$ is divided by $DD'$, the remainder is:


 $\textbf{(A)}\ R+R'D\qquad \textbf{(B)}\ R'+RD\qquad \textbf{(C)}\ RR'\qquad \textbf{(D)}\ R\qquad \textbf{(E)}\ R'$

 \textbf{Answer: }A



\textbf{Solution}

We are given $P = QD + R$ and $Q = D'Q' + R'$.

 Plugging the second equation into the first yields:

 $P = (D'Q' + R')D + R$
$P = (DD')Q' + (R'D + R)$

 If we divide $P$ by $DD'$, the quotient would be $Q'$, and the remainder would be $R'D + R$, which is option $\fbox{A}$.

 


	\item (\textit{AHSME 1967}) If $x$ is real and positive and grows beyond all bounds, then $\log_3{(6x-5)}-\log_3{(2x+1)}$ approaches:


 $\textbf{(A)}\ 0\qquad \textbf{(B)}\ 1\qquad \textbf{(C)}\ 3\qquad \textbf{(D)}\ 4\qquad \textbf{(E)}\ \text{no finite number}$

 \textbf{Answer: }B



\textbf{Solution}

Since $\log_b x - \log_b y = \log_b \frac{x}{y}$, the expression is equal to $\log_3 \frac{6x - 5}{2x + 1}$.

 The expression $\frac{6x - 5}{2x + 1}$ is equal to $3 - \frac{8}{2x + 1}$.  As $x$ gets large, the second term approaches $0$, and thus $\frac{6x - 5}{2x + 1}$ approaches $3$.  Thus, the expression approaches $\log_3 3$, which is $1$.

 Alternately, we divide the numerator and denominator of $\frac{6x - 5}{2x + 1}$ by $x$ to get $\frac{6 - \frac{5}{x}}{2 + \frac{1}{x}}$.  As $x$ grows large, both fractions approach $0$, leaving $\frac{6}{2} = 3$, and so the expression approaches $\log_3 3 = 1$.

 With either reasoning, the answer is $\fbox{B}$



	\item (\textit{AHSME 1967}) The number of solution-pairs in the positive integers of the equation $3x+5y=501$ is:


 $\textbf{(A)}\ 33\qquad \textbf{(B)}\ 34\qquad \textbf{(C)}\ 35\qquad \textbf{(D)}\ 100\qquad \textbf{(E)}\ \text{none of these}$

 \textbf{Answer: }A



\textbf{Solution}

We have $y = \frac{501 - 3x}{5}$.  Thus, $501 - 3x$ must be a positive multiple of $5$.  If $x = 2$, we find our first positive multiple of $5$.  From there, we note that $x = 2 + 5k$ will always return a multiple of $5$ for $501 - 3x$.  Our first solution happens at $k=0$.  

 
We now want to find the smallest multiple of $5$ that will work.  If $x = 2 + 5k$, then we have $501 - 3x = 501 - 3(2 + 5k)$, or $495 - 15k$.  When $k = 32$, the expression is equal to $15$, and when $k = 33$, the expression is equal to $0$, which will no longer work.

 
Thus, all integers from $k = 0$ to $k = 32$ will generate an $x = 2 + 5k$ that will be a positive integer, and which will in turn generate a $y$ that is also a positive integer.  So, the answer is $\fbox{A}$.

 


	\item (\textit{AHSME 1967}) For every odd number $p>1$ we have:


 $\textbf{(A)}\ (p-1)^{\frac{1}{2}(p-1)}-1 \; \text{is divisible by} \; p-2\qquad \textbf{(B)}\ (p-1)^{\frac{1}{2}(p-1)}+1 \; \text{is divisible by} \; p\\ \textbf{(C)}\ (p-1)^{\frac{1}{2}(p-1)} \; \text{is divisible by} \; p\qquad \textbf{(D)}\ (p-1)^{\frac{1}{2}(p-1)}+1 \; \text{is divisible by} \; p+1\\ \textbf{(E)}\ (p-1)^{\frac{1}{2}(p-1)}-1 \; \text{is divisible by} \; p-1$

 \textbf{Answer: }A



\textbf{Solution}

Given that $p$ is odd, $p-1$ must be even, therefore ${{\frac{1}{2}}(p-1)}$ must be an integer, which will be denoted as n.
 \[(p-1)^n-1\]By sum and difference of powers
 \[=((p-1)-1)((p-1)^{n-1}+\cdots+1^{n-1})\]
 \[=(p-2)((p-1)^{n-1}+\cdots+1^{n-1})\]$p-2$ divide $(p-1)^{{\frac{1}{2}}(p-1)}$$\fbox{A}$.

 


	\item (\textit{AHSME 1967}) If one uses only the tabular information $10^3=1000$, $10^4=10,000$, $2^{10}=1024$, $2^{11}=2048$, $2^{12}=4096$, $2^{13}=8192$, then the strongest statement one can make for $\log_{10}{2}$ is that it lies between:


 $\textbf{(A)}\ \frac{3}{10} \; \text{and} \; \frac{4}{11}\qquad \textbf{(B)}\ \frac{3}{10} \; \text{and} \; \frac{4}{12}\qquad \textbf{(C)}\ \frac{3}{10} \; \text{and} \; \frac{4}{13}\qquad \textbf{(D)}\ \frac{3}{10} \; \text{and} \; \frac{40}{132}\qquad \textbf{(E)}\ \frac{3}{11} \; \text{and} \; \frac{40}{132}$

 \textbf{Answer: }C



\textbf{Solution}

Since $1024$ is greater than $1000$.

 $\log (1024) > 3$
$10 * \log (2) > 3$

 and $\log (2) > 3/10$.

 
Similarly, $8192 < 10000$, so $\log (8192) < 4$
$13 * \log (2) < 4$

 and $\log (2) < 4/13$

 
Therefore  $3/10 < \log 2 < 4/13$ so the answer is $\fbox{C}$



	\item (\textit{AHSME 1967}) Two candles of the same length are made of different materials so that one burns out completely at a uniform rate in $3$ hours and the other in $4$ hours.  At what time P.M. should the candles be lighted so that, at 4 P.M., one stub is twice the length of the other?


 $\textbf{(A)}\ 1:24\qquad \textbf{(B)}\ 1:28\qquad \textbf{(C)}\ 1:36\qquad \textbf{(D)}\ 1:40\qquad \textbf{(E)}\ 1:48$

 \textbf{Answer: }C



\textbf{Solution}

If the candles both have length $\ell$, then the candle that burns in $3$ hours has a stub of $\ell$ at $0$ minutes, and a stub of $0$ at $180$ minutes.  Since the candle burns at a constant rate (i.e. linearly), the stub length of this candle $t$ minutes after being lit is $f(t) = \frac{\ell}{180}(180 - t)$, since $f(0) = \ell$ and $f(180) = 0$.

 Similarly, for the second candle that burns out in $240$ minutes, $g(t) = \frac{\ell}{240}(240 - t)$

 Since the first candle burns out faster, if the two candles are lighted simultaneously, it will always have a shorter stub.  The problem asks for when $g(t) = 2f(t)$.  Solving this equation gives:

 $\frac{\ell}{240}(240 - t) = 2\frac{\ell}{180}(180 - t)$
$240 - t = \frac{480}{180}(180 - t)$
$240 - t = 480 - \frac{480}{180}t$
$\frac{8}{3} t - t = 480 - 240$
$t = \frac{3}{5} \cdot 240$
$t = 144$

 So, the second candle will have a stub twice as big as the first candle $144$ minutes after they are both lit.  If we want this to happen at $4$ PM, the candles have to be lit $144$ minutes earlier, or $2$ hours and $24$ minutes earlier.  This is at $\text{1:36 PM}$, which is option $\fbox{C}$



	\item (\textit{AHSME 1967}) Given the two hypotheses: $\text{I}$ Some Mems are not Ens and $\text{II}$ No Ens are Veens.  If "some" means "at least one," we can conclude that:


 $\textbf{(A)}\ \text{Some Mems are not Veens}\qquad \textbf{(B)}\ \text{Some Vees are not Mems}\\ \textbf{(C)}\ \text{No Mem is a Vee}\qquad \textbf{(D)}\ \text{Some Mems are Vees}\\ \textbf{(E)}\ \text{Neither} \; \textbf{(A)} \; \text{nor} \; \textbf{(B)} \; \text{nor} \; \textbf{(C)} \; \text{nor} \; \textbf{(D)} \; \text{is deducible from the given statements}$

 \textbf{Answer: }E



\textbf{Solution}

$\fbox{E}$



	\item (\textit{AHSME 1967}) $\overline{AB}$ is a diameter of a circle.  Tangents $\overline{AD}$ and $\overline{BC}$ are drawn so that $\overline{AC}$ and $\overline{BD}$ intersect in a point on the circle.  If $\overline{AD}=a$ and $\overline{BC}=b$, $a \not= b$, the diameter of the circle is:


 $\textbf{(A)}\ |a-b|\qquad \textbf{(B)}\ \frac{1}{2}(a+b)\qquad \textbf{(C)}\ \sqrt{ab} \qquad \textbf{(D)}\ \frac{ab}{a+b}\qquad \textbf{(E)}\ \frac{1}{2}\frac{ab}{a+b}$

 \textbf{Answer: }C



\textbf{Solution}

$\fbox{C}$



	\item (\textit{AHSME 1967}) A dealer bought $n$ radios for $d$ dollars, $d$ a positive integer.  He contributed two radios to a community bazaar at half their cost.  The rest he sold at a profit of $8 on each radio sold.  If the overall profit was $72, then the least possible value of $n$ for the given information is:


 $\textbf{(A)}\ 18\qquad \textbf{(B)}\ 16\qquad \textbf{(C)}\ 15\qquad \textbf{(D)}\ 12\qquad \textbf{(E)}\ 11$

 \textbf{Answer: }D



\textbf{Solution}

The first $2$ radios are sold for $\frac{d}{2}$ dollars each, for income of $2 \cdot \frac{d}{2}$, or simply $d$.

 The remaining $n-2$ radios are sold for $d+8$ dollars each, for income of $(n-2)(d+8)$, which expands to $nd +8n -2d - 16$.

 The amount invested in buying $n$ radios for $d$ dollars is $nd$.

 Thus, the total profit is $(d) + (nd - 2d + 8n - 16) - nd$, which simplifies to $8n - d - 16$.

 Setting this equal to $72$ gives $8n - d - 16 = 72$, or $n = 11 + \frac{d}{8}$.

 Since $n$ must be an integer, $d$ must be a multiple of $8$.  If we want to minimize $n$, we want to minimize $d$.  Since $d>0$, setting $d=8$ yields $n=12$, giving the answer of $\fbox{D}$



	\item (\textit{AHSME 1967}) Let $D=a^2+b^2+c^2$, where $a$, $b$, are consecutive integers and $c=ab$.  Then $\sqrt{D}$ is:


 $\textbf{(A)}\ \text{always an even integer}\qquad \textbf{(B)}\ \text{sometimes an odd integer, sometimes not}\\ \textbf{(C)}\ \text{always an odd integer}\qquad \textbf{(D)}\ \text{sometimes rational, sometimes not}\\ \textbf{(E)}\ \text{always irrational}$

 \textbf{Answer: }C



\textbf{Solution}

Let $a=x, b=x+1, c = x(x+1)$.  Then $D = x^2 + (x+1)^2 + x^2(x+1)^2$, which simplifies to $x^4 + 2x^3 + 3x^2 + 2x + 1$.

 From the options, we want to test if $D$ is always a perfect square.  Because the polynomial expression for $D$ is quartic, if it is a perfect square, it would be the square of a quadratic expression.  Thus, $D$ could be written in the form $(Ax^2 + Bx + C)^2$ for some $(A, B, C)$.

 Setting $(Ax^2 + Bx + C)^2 = x^4 + 2x^3 + 3x^2 + 2x + 1$, we can compare coefficients.  From the $x^4$ coefficient, we get $A = \pm 1$.  Note that if $(Ax^2 + Bx^2 + C)^2$ works, so does $(-Ax^2 - Bx - C)^2$, so we can arbitrarily pick $A=1$.

 We now have $(x^2 + Bx + C)^2= x^4 + 2x^3 + 3x^2 + 2x + 1$.  Setting the cubic terms equal gives $2B = 2$, or $B = 1$.  This leaves $(x^2 + x + C)^2 =  x^4 + 2x^3 + 3x^2 + 2x + 1$.  We can quickly inspect the constant term to determine that $C = \pm 1$.  We reject $C = -1$, since the quadratic and linear terms won't match up, which leaves $(x^2 + x + 1)^2$ as the only possibility - and, in fact, it works.

 Thus, $D$ is always the square of an integer - namely $x^2 + x + 1$. This in turn means that $\sqrt{D}$ is always rational, which leaves choices $A, B, C$ as the only possible correct answers.

 The question now is whether $\sqrt{D}$, or $x^2 + x + 1$, is odd, even, or could be both.  We have two cases for $x$:

 If $x \equiv 0 \pmod 2$, then $x^2 \equiv 0^2 \pmod 2$.  This means $x^2 + x + 1 \equiv 0 + 0 + 1 \equiv 1 \pmod 2$.  

 If $x \equiv 1 \pmod 2$, then $x^2 \equiv 1^2 \pmod 2$, and $x^2 + x + 1 \equiv 1 + 1 + 1 \equiv 1 \pmod 2$.  

 Either way, $x^2 + x + 1$ is an odd integer, and the answer is $\fbox{C}$



	\item (\textit{AHSME 1967}) In quadrilateral $ABCD$ with diagonals $\overline{AC}$ and $\overline{BD}$ intersecting at $O$, $\overline{BO}=4$, $\overline{OD}=6$, $\overline{AO}=8$, $\overline{OC}=3$, and $\overline{AB}=6$.  The length of $\overline{AD}$ is:


 $\textbf{(A)}\ 9\qquad \textbf{(B)}\ 10\qquad \textbf{(C)}\ 6\sqrt{3}\qquad \textbf{(D)}\ 8\sqrt{2}\qquad \textbf{(E)}\ \sqrt{166}$

 \textbf{Answer: }E



\textbf{Solution 1}

After drawing the diagram, we see that we actually have a lot of lengths to work with. Considering triangle ABD, we know values of $AB, BD(BD = BO + OD)$, but we want to find the value of $AD$. We can apply stewart's theorem now, letting $m = 4, n = 6, AD = X, AB = 6$, and we have $10 \cdot 6 \cdot 4 + 8 \cdot 8 \cdot 10 = x^2 + 36 \cdot 6$, and we see that $x = \boxed{\textbf{(E)}~\sqrt{166}}$




\textbf{Solution 2}


                    \begin{figure}[H]
                        \centering
                        \includegraphics[width=0.5\textwidth]{Solution/AHSME_I_1967_32_Solution_Figure_1.png}
                    \end{figure}
                    (Diagram not to scale)

 Since $AO \cdot OC = BO \cdot OD$, $ABCD$ is cyclic through power of a point. From the given information, we see that $\triangle{AOB}\sim \triangle{DOC}$ and $\triangle{BOC} \sim \triangle{AOD}$. Hence, we can find $CD=\frac{9}{2}$ and $AD=2 \cdot BC$. Letting $BC$ be $x$, we can use Ptolemy's to get
 \[6 \cdot \frac{9}{2} + 2x^2=10 \cdot 11 \implies x=\sqrt{\frac{83}{2}}\]Since we are solving for $AD=2x=2\cdot\sqrt{\frac{83}{2}}=\sqrt{4\cdot\frac{83}{2}} = \boxed{\textbf{(E)}~\sqrt{166}}$

 - PhunsukhWangdu

 



\textbf{Solution 3:}

The solution says it all. Since $\angle AOD$ is supplementary to $\angle AOB$, $cos(\angle AOD) = cos(180^{\circ} - \angle AOB)=-cos(\angle AOB)$. The law of cosines on $\triangle AOB$ gives us $cos(\angle AOB)=\frac {8^2+4^2-6^2}{(2)(8)(4)}=\frac {11}{16}$. Again, we can use the law of cosines on $\triangle AOD$, which gives us 
 \[AD=\sqrt {8^2+6^2-2(8)(6)cos(\angle AOD)}\]
 \[=\sqrt {8^2+6^2-(2)(8)(6)cos(\angle 180^{\circ} - AOB)}\]
 \[=\sqrt {8^2+6^2+(2)(8)(6)cos(\angle AOB)}\]
 \[=\sqrt {8^2+6^2+(2)(8)(6)(\frac {11}{16})}\]
 \[=\sqrt{166}\]which gives us $\boxed{\textbf{(E)}~\sqrt{166}}$

 Note that this solution works even if the quadrilateral is not cyclic, and in general, it works if an angle's supplement is known. (Anyone come from aops volume 2 lmao.)

 -Tanstorm

 


	\item (\textit{AHSME 1967}) 
                    \begin{figure}[H]
                        \centering
                        \includegraphics[width=0.5\textwidth]{Problem/AHSME_I_1967_33_Problem_Figure_1.png}
                    \end{figure}
                    


 In this diagram semi-circles are constructed on diameters $\overline{AB}$, $\overline{AC}$, and $\overline{CB}$, so that they are mutually tangent.  If $\overline{CD} \bot \overline{AB}$, then the ratio of the shaded area to the area of a circle with $\overline{CD}$ as radius is:


 $\textbf{(A)}\ 1:2\qquad \textbf{(B)}\ 1:3\qquad \textbf{(C)}\ \sqrt{3}:7\qquad \textbf{(D)}\ 1:4\qquad \textbf{(E)}\ \sqrt{2}:6$

 \textbf{Answer: }D



\textbf{Solution}

To make the problem much simpler while staying in the constraints of the problem, position point $C$ halfway between $A$ and $B$. Then, call $\overline{AC} = \overline{BC}=r$ . The area of the shaded region is then 
 \[\frac{ \pi r^2 - \pi (r/2)^2 - \pi (r/2)^2}{2}=\frac{\pi r^2}{4}\]Because $\overline{CD}=r$ the area of the circle with $\overline{CD}$ as radius is $\pi r^2$.Our ratio is then 
 \[\frac{\pi r^2}{4} : \pi r^2 = 1:4\]

 Which corresponds with answer $\fbox{D}$



	\item (\textit{AHSME 1967}) Points $D$, $E$, $F$ are taken respectively on sides $AB$, $BC$, and $CA$ of triangle $ABC$ so that $AD:DB=BE:CE=CF:FA=1:n$.  The ratio of the area of triangle $DEF$ to that of triangle $ABC$ is:


 $\textbf{(A)}\ \frac{n^2-n+1}{(n+1)^2}\qquad \textbf{(B)}\ \frac{1}{(n+1)^2}\qquad \textbf{(C)}\ \frac{2n^2}{(n+1)^2}\qquad \textbf{(D)}\ \frac{n^2}{(n+1)^2}\qquad \textbf{(E)}\ \frac{n(n-1)}{n+1}$

 \textbf{Answer: }A



\textbf{Solution}

WLOG, let's assume that $\triangle ABC$ is equilateral. Therefore, $[ABC]=\frac{(1+n)^2\sqrt3}{4}$ and $[DBE]=[ADF]=[EFC]=n \cdot \sin(60)/2$. Then $[DEF]=\frac{(n^2-n+1)\sqrt3}{4}$. Finding the ratio yields $\fbox{A}$.  -Dark\_Lord

 


	\item (\textit{AHSME 1967}) The roots of $64x^3-144x^2+92x-15=0$ are in arithmetic progression.  The difference between the largest and smallest roots is:


 $\textbf{(A)}\ 2\qquad \textbf{(B)}\ 1\qquad \textbf{(C)}\ \frac{1}{2}\qquad \textbf{(D)}\ \frac{3}{8}\qquad \textbf{(E)}\ \frac{1}{4}$

 \textbf{Answer: }B



\textbf{Solution 1}

By \href{https://artofproblemsolving.com/wiki/index.php?title=Vieta\%27s\_Formulas}{Vieta's Formulas}, the sum of the roots is $-\frac{-144}{64} = \frac{9}{4}$.  Because the roots are in arithmetic progression, the middle root is the average of the other two roots, and is also the average of all three roots.  Therefore, $\frac{\frac{9}{4}}{3} = \frac{3}{4}$ is the middle root.

 The other two roots have a sum of $\frac{9}{4} - \frac{3}{4} = \frac{3}{2}$.  By Vieta on the original cubic, the product of all $3$ roots is $-\frac{-15}{64} = \frac{15}{64}$, so the product of the remaining two roots is $\frac{\frac{15}{64}}{\frac{3}{4}} = \frac{5}{16}$.  If the sum of the two remaining roots is $\frac{3}{2} = \frac{24}{16}$, and the product is $\frac{5}{16}$, the two remaining roots are also the two roots of $16x^2 - 24x + 5 = 0$.

 The two remaining roots are thus $x = \frac{-B \pm \sqrt{B^2 - 4AC}}{2A}$, and they have a difference of $\frac{2\sqrt{B^2 - 4AC}}{2A}$.  Plugging in gives $\frac{\sqrt{24^2 - 4(16)(5)}}{16}$, which is equal to $1$, which is answer $\fbox{B}$.

 



\textbf{Solution 2}

Because the roots are in arithmetic progression, let the roots be $r-d, r,$ and $r+d$ with $d \geq 0$. From \href{https://artofproblemsolving.com/wiki/index.php?title=Vieta\%27s\_Formulas}{Vieta's Formulas}, we know that $(r-d)+r+(r+d)=\frac{144}{64}$, so $3r=\frac 9 4$, and $r=\frac 3 4$.

 By using Vieta again, we see that $(r-d)(r)(r+d)=\frac{15}{64}$. Substituting $r=\frac 3 4$ into this equation, we can now solve for $d$:\begin{align*}(\frac 3 4-d)(\frac 3 4)(\frac 3 4+d) &= \frac{15}{64} \\\frac 3 4(\frac 9{16}-d^2) &= \frac{15}{64} \\\frac 9{16}-d^2 &= \frac 5{16} \\d^2 &= \frac 4{16} = \frac 1 4 \\d &= \pm \frac 1 2\end{align*}
 Because we defined $d \geq 0$, $d=\frac 1 2$. Therefore, the difference between the largest and smallest roots, $r+d$ and $r-d$, is $2d=2 \cdot \frac 1 2 = \boxed{\textbf{(B) }1}$.

 


	\item (\textit{AHSME 1967}) Given a geometric progression of five terms, each a positive integer less than $100$.  The sum of the five terms is $211$.  If $S$ is the sum of those terms in the progression which are squares of integers, then $S$ is:


 $\textbf{(A)}\ 0\qquad \textbf{(B)}\ 91\qquad \textbf{(C)}\ 133\qquad \textbf{(D)}\ 195\qquad \textbf{(E)}\ 211$

 \textbf{Answer: }C



\textbf{Solution}

Let the first term be $a$ and the common ratio be $r$, and WLOG let $r \ge 1$.  The five terms are $a, ar, ar^2, ar^3, ar^4$, and the sum is $a(1 + r + r^2 + r^3 + r^4)$.  Clearly $r$ must be rational for all terms to be integers.  If $r$ were an integer, it could not be $1$, since $a$ would equal $\frac{211}{1 + r + r^2 + r^3 + r^4}$, which is not an integer.  In fact, quickly testing $r = 2, 3, 4$ shows that $r$ cannot be an integer.

 We now consider non-integers.  If $r = \frac{x}{y}$ and $\gcd(x, y) = 1$, then $a$ would have to be divisible by $y^4$, since $ar^4 = a\frac{x^4}{y^4}$ is an integer.  If $y = 3$, then $a$ would have to be a multiple of $81$, which would make the five terms sum to at least $5 \cdot 81$.  It only gets worse if $y > 3$.  Thus, $y = 2$, and $a$ is a multiple of $2^4 = 16$.  Let $a = 16k$.

 We now look at the last term, $ar^4 = 16k\frac{x^4}{16} = kx^4$.  The smallest allowable values for $x$, given that $x$ cannot be even, are $3$ and $5$.  If $x = 5$, the last term will be way too big.  Thus, $x = 3$ is the only possibility.

 We now have a sum of $a(1 + r + r^2 + r^3 + r^4) = 211$, and we know that $r = \frac{3}{2}$ is the only possibility.  The terms in the parentheses happen to equal $\frac{211}{16}$ when you plug them in, so $a = 16$.

 Thus, the terms are $16, 24, 36, 54, 81$, and the first, third, and fifth terms are squares, with a sum of $133$, which is answer $\fbox{C}$.

 


	\item (\textit{AHSME 1967}) Segments $AD=10$, $BE=6$, $CF=24$ are drawn from the vertices of triangle $ABC$, each perpendicular to a straight line $RS$, not intersecting the triangle.  Points $D$, $E$, $F$ are the intersection points of $RS$ with the perpendiculars.  If $x$ is the length of the perpendicular segment $GH$ drawn to $RS$ from the intersection point $G$ of the medians of the triangle, then $x$ is:


 $\textbf{(A)}\ \frac{40}{3}\qquad \textbf{(B)}\ 16\qquad \textbf{(C)}\ \frac{56}{3}\qquad \textbf{(D)}\ \frac{80}{3}\qquad \textbf{(E)}\ \text{undetermined}$

 \textbf{Answer: }A



\textbf{Solution}

$\fbox{A}$

 WLOG let $RS$ be the x-axis, or at least horizontal. The three lengths represent the y-coordinates of points $A,B,C$. As $G$ is by definition the average of $A,B,C$, it's coordinates are the average of the coordinates of $A,B,$ and $C$. Hence, the y-coordinate of $G$, which is also the distance from $G$ to $RS$, is the average of the y-coordinates of $A,B,$ and $C$, or $\frac{10+6+24}{3}$, hence the result.

 


	\item (\textit{AHSME 1967}) Given a set $S$ consisting of two undefined elements "pib" and "maa", and the four postulates: $P_1$: Every pib is a collection of maas, $P_2$: Any two distinct pibs have one and only one maa in common, $P_3$: Every maa belongs to two and only two pibs, $P_4$: There are exactly four pibs.  Consider the three theorems: $T_1$: There are exactly six maas, $T_2$: There are exactly three maas in each pib, $T_3$: For each maa there is exactly one other maa not in the same pib with it.  The theorems which are deducible from the postulates are:


 $\textbf{(A)}\ T_3 \; \text{only}\qquad \textbf{(B)}\ T_2 \; \text{and} \; T_3 \; \text{only} \qquad \textbf{(C)}\ T_1 \; \text{and} \; T_2 \; \text{only}\\ \textbf{(D)}\ T_1 \; \text{and} \; T_3 \; \text{only}\qquad \textbf{(E)}\ \text{all}$

 \textbf{Answer: }E



\textbf{Solution}

$\fbox{E}$



	\item (\textit{AHSME 1967}) Given the sets of consecutive integers $\{1\}$,$\{2, 3\}$,$\{4,5,6\}$,$\{7,8,9,10\}$,$\; \cdots \;$, where each set contains one more element than the preceding one, and where the first element of each set is one more than the last element of the preceding set.  Let $S_n$ be the sum of the elements in the nth set.  Then $S_{21}$ equals:


 $\textbf{(A)}\ 1113\qquad \textbf{(B)}\ 4641 \qquad \textbf{(C)}\ 5082\qquad \textbf{(D)}\ 53361\qquad \textbf{(E)}\ \text{none of these}$

 \textbf{Answer: }B



\textbf{Solution}

The last element of the 21st set is $\frac{21\times 22}{2}=231$, and hence the first element of the 21st set is $211$.So $S_{21}=\frac{231\times 232}{2}-\frac{210\times 211}{2}=4641$, hence our answer is $\fbox{B}$



	\item (\textit{AHSME 1967}) Located inside equilateral triangle $ABC$ is a point $P$ such that $PA=8$, $PB=6$, and $PC=10$.  To the nearest integer the area of triangle $ABC$ is:


 $\textbf{(A)}\ 159\qquad \textbf{(B)}\ 131\qquad \textbf{(C)}\ 95\qquad \textbf{(D)}\ 79\qquad \textbf{(E)}\ 50$

 \textbf{Answer: }D



\textbf{Solution 1}


                    \begin{figure}[H]
                        \centering
                        \includegraphics[width=0.5\textwidth]{Solution/AHSME_I_1967_40_Solution_Figure_1.png}
                    \end{figure}
                    

 Notice that $6^2+8^2=10^2.$ That makes us want to construct a right triangle.

 Rotate $\triangle APC$ $60^{\circ}$ about A. Note that $\triangle PAC \cong \triangle P'AB$, so 
 \[\angle P'AP = \angle PAB + \angle P'AB = \angle PAB + \angle PAC = 60^{\circ}.\]

 Therefore, $\triangle APP'$ is equilateral, so $P'P=8$, which means $\angle P'PB = 90^{\circ}.$ 

 Let $\angle BP'P = \alpha .$ Notice that $\cos\alpha = \frac{8}{10}=\frac{4}{5},$ and $\sin\alpha = \frac{3}{5}.$

 Applying the Law of Cosines to $\triangle APC$ (remembering $\angle APC = \angle AP'B$):
 \begin{align*} AC^2 &= 10^2+8^2-2\cdot10\cdot8\cdot \cos(60^{\circ}+\alpha)\\&= 164-160(\cos60\cos\alpha-\sin60\sin\alpha)\\&= 164-160(\frac{2}{5}-\frac{3\sqrt3}{10}) \\&= 164-16(4-3\sqrt3) \\ &= 100+48\sqrt3.\end{align*}


 We want to find the area of $\triangle ABC$, which is 
 \[AC^2\cdot\frac{\sqrt3}{4}=25\sqrt3+36\approx\boxed{(D) 79}.\]




\textbf{Solution 2 (Magic Formula)}

Fun formula:  Given a point whose distances from the vertices of an equilateral triangle are $a$, $b$, and $c$, the side length of the triangle is:

 
 \[s=\sqrt{\frac{1}{2}\left(a^2+b^2+c^2\pm\sqrt{6(a^2b^2+b^2c^2+c^2a^2)-3(a^4+b^4+c^4)}\right)}\]

 Given that the area of an equilateral triangle is $\frac{\sqrt{3}}{4}s^2$, the answer is:

 
 \begin{align*} [ABC] &= \frac{\sqrt{3}}{4}\cdot\frac{1}{2}\left(a^2+b^2+c^2\pm\sqrt{6(a^2b^2+b^2c^2+c^2a^2)-3(a^4+b^4+c^4)}\right)\\ &= \frac{\sqrt{3}}{8}\left(200\pm\sqrt{6\cdot 16(3^2 4^2+4^2 5^2+5^2 3^2)-3\cdot16(3^4+4^4+5^4)}\right)\\ &= \frac{\sqrt{3}}{8}\left(200\pm\sqrt{96(144+400+225)-48(81+256+625)}\right)\\ &= \frac{\sqrt{3}}{8}\left(200\pm\sqrt{96\cdot769-48\cdot962}\right) = \frac{\sqrt{3}}{8}\left(200\pm\sqrt{96\cdot769-96\cdot481}\right)\\ &= \frac{\sqrt{3}}{8}\left(200\pm\sqrt{96\cdot288}\right) = \frac{\sqrt{3}}{8}\left(200\pm\sqrt{96\cdot96\cdot3}\right)\\ &= 25\sqrt{3}\pm36 \approx \{6.5, \text{or } 78.5\} \end{align*}

$6.5$ is not a choice, therefore the answer is $\boxed{\textbf{(D) }79}$.

 (Note that the $6.5$ answer is actually the solution for when point $P$ is \textit{exterior} to $\triangle ABC$.)

 



\textbf{Solution 3}

Rotate $P$ and $B$ $60^{\circ}$ CCW around $A$, becoming $X$ and $C$.  Rotate $P$ and $C$ $60^{\circ}$ CCW around $B$, becoming $Y$ and $A$.  Rotate $P$ and $A$ $60^{\circ}$ CCW around $C$, becoming $Z$ and $B$:

 
                    \begin{figure}[H]
                        \centering
                        \includegraphics[width=0.5\textwidth]{Solution/AHSME_I_1967_40_Solution_Figure_2.png}
                    \end{figure}
                    

 Notice that since $\triangle AXC\cong\triangle APB$, $\triangle BYA\cong\triangle BPC$, and $\triangle CZB\cong\triangle CPA$, then

 
 \[[AYBZCX]=2\cdot[ABC]\]

 Now the area of the big hexagon is easy to compute since it's comprised of 3 equilateral triangle and 3 right triangles:

 
 \begin{align*}[ABC] &= \frac{1}{2}[AYBZCX] = \frac{1}{2}\left(\underbrace{3\cdot\frac12\cdot6\cdot8}_{\text{3 right triangles}}+\underbrace{\frac{\sqrt{3}}{4}\left(6^2+8^2+10^2\right)}_{\text{3 equilateral triangles}}\right)\\ &= \frac{1}{2}\left(72+\frac{\sqrt{3}}{4}\cdot200)\right) = 36+25\sqrt{3}\\ &\approx \boxed{\textbf{(D) }79} \end{align*}





\textbf{Solution 4(Answer Choices, Approximation)}

Let $s$ be the side length of $ABC.$ Notice that $s\le 14$ by the triangle inequality. This means that
 \[[ABC]\le\dfrac{14^2\sqrt{3}}{2}\approx 84.87.\] This automatically rules out choices $A, B,$ and $C.$ Now, we will look at if the area is $50$. By the equilateral triangle area formula, $s$ would equal $10\sqrt{\dfrac{2}{\sqrt{3}}}\approx 10.75.$ This is very close to $10.$ If $s=10,$ $\angle APB=90$ and $\angle APC, \angle BPC<90$ by the Pythagorean theorem and Pythagorean inequalities. Thus, 
 \[\angle APB+\angle APC+\angle BPC<270.\]$\angle APB+\angle APC+\angle BPC$ needs to be $360,$ and it probably cannot increase by more than $90$ by just adding $0.75$ to $s$ (more rigorous proof below) Thus, the only viable answer choice is $\boxed{79}.$

            \begin{itemize}
        \item In fact, for $s<\sqrt{136}\approx 11.66$ ($[ABC]<34\sqrt{3}\approx58.88$) $\angle APC, \angle BPC<90$ still holds. As $\angle APB <180,$ 
 \[\angle APB+\angle APC+\angle BPC<360.\] Thus, we know for sure that the nearest integer to the area cannot be $50.$
            \end{itemize}
            


\end{enumerate}